% Generated by tools/build_docs.py from docs/templates/memo.tex.tmpl. Do not edit memo.tex by hand.
% Compile with xelatex. Independent analysis using public information and synthetic data. Not affiliated with or endorsed by Varaha.
\documentclass[10pt]{article}
\usepackage[a4paper,margin=1.35cm,top=1.1cm,bottom=1.0cm]{geometry}
\usepackage{fontspec}
\setmainfont{Inter}[Scale=0.92]
\newfontfamily\rupeefont{DejaVu Sans}[Scale=0.88]
\newcommand{\rupee}{{\rupeefont ₹}}
\usepackage{xcolor,tabularx,booktabs,enumitem,array}
\definecolor{accent}{HTML}{0F766E}\definecolor{ink}{HTML}{111827}\definecolor{muted}{HTML}{6B7280}
\color{ink}
\setlength{\parindent}{0pt}\setlength{\parskip}{3pt}
\setlist{leftmargin=1.1em,itemsep=1pt,topsep=1pt,parsep=0pt}
\newcommand{\hd}[1]{\vspace{4pt}{\color{accent}\bfseries\large #1}\par\vspace{1pt}}
\renewcommand{\arraystretch}{1.1}
\newcolumntype{Y}{>{\raggedright\arraybackslash}X}\newcolumntype{L}[1]{>{\raggedright\arraybackslash}p{#1}}
\pagestyle{empty}
\begin{document}
{\small\color{muted} \textbf{To:} COO \quad \textbf{From:} Operations strategy (Founder's Office) \quad \textbf{Date:} December 2026 \quad \textbf{Scope:} a 450-kiln distributed network modelled on Agnisparsh}\\[-2pt]
{\color{accent}\rule{\linewidth}{1.2pt}}\\[2pt]
{\LARGE\bfseries Delivering 100k t from distributed biochar in 2027: three risks and the Q1 plan}\par\vspace{4pt}

\hd{Bottom line}
\textbf{The network reaches 100k t only at its design point, and even there our P50 is 68.4k t (P10–P90 52.7k–87.9k; P(≥100k) 2.1\%), because almost every nominal input sits near the favourable end of its range.} Q1 has to buy about 32k t of contingency. The cheapest tonnes are the ones we already produce and then lose at dMRV, at \rupee{}1,336/t to recover against about \rupee{}9,728/t for new capacity.

\hd{What has to be true}
\begin{itemize}
\item \textbf{Throughput of 4.5 t dry feedstock per kiln-day:} 2.5 t batches (a 3\,m³-class kiln) at 1.8 batches/day, the single-shift cycle ceiling. On published Kon-Tiki values throughput is 2.25 t/kiln-day and the network delivers 44.1k t. This is the first number to measure.
\item \textbf{Feedstock dried to ≤15\% moisture.} The CH₄ deduction then falls from 41\% to 7\% of gross removals.
\item \textbf{About 86\% of batches pass dMRV.} At the design point 26\% of gross removals is lost before credit, mostly to rejections.
\item \textbf{Jan–Jun carries 53\% of the year and the monsoon only 13\%.} Slippage before June cannot be recovered. In our synthetic year, actuals ran at 80\% of plan from January.
\end{itemize}

\hd{Three delivery risks (quantified on a synthetic operating year)}
{\small
\begin{tabularx}{\linewidth}{@{}L{3.7cm}YrL{2.3cm}@{}}
\toprule
\textbf{Risk} & \textbf{Evidence} & \textbf{t at risk} & \textbf{Likelihood / impact}\\\midrule
Kiln throughput below design & 1.66 batches per running kiln-day vs 1.8 plan; utilisation 89\%; idle kiln-days from stockouts, vacancies and moves & 17.9k & High / High\\
dMRV rejections & Acceptance 84.6\% vs 86\% target; top two reasons (missing temperature log, missing / unusable photos) are 44\% of rejections; new operators 76\% vs 85\% & 14.4k & High / High\\
Biomass supply and payment integrity & Paid within SLA: 88\% network, 55\% worst site. At the worst site, supply fell 4 weeks after payments broke SLA (1.04× → 0.52× peer level). This is a buyer and rating risk as well as a tonne risk & 4.4k & Medium / High\\\bottomrule
\end{tabularx}}\\[1pt]
{\footnotesize\color{muted}Calibration risk outside ops control: persistence factor and CH₄ treatment under the methodology (up to 32.3k t downside).}

\hd{Five Q1 actions}
{\small
\begin{tabularx}{\linewidth}{@{}L{0.3cm}YL{2.5cm}L{3.6cm}L{1.9cm}@{}}
\toprule
& \textbf{Action} & \textbf{Owner} & \textbf{Metric (target)} & \textbf{Date}\\\midrule
1 & Stand up a dMRV QA desk: daily evidence review, thermocouple and scale checks, a 14-day buddy for new lead operators & MRV lead & Acceptance ≥ 90\% (now 85\%) & 15 Feb 2027\\
2 & Weigh every batch and time 5 sites for two weeks, then choose the capacity lever: second shift on 25\% of kilns (+15.6k t) or 45 more kilns (+8.2k t) & Head of field operations & t dry per running kiln-day vs 4.5 design & Baseline 31 Jan; decision 31 Mar\\
3 & 14-day village payment SLA with a weekly breach list; escalate any site below 75\% & Finance lead + community lead & ≥ 90\% paid on time at every site (worst now 55\%) & Weekly from 15 Jan\\
4 & Moisture gate: a meter at every site, no load above 15\%, covered drying yards & Site supervisors (QA) & ≥ 90\% of batches ≤ 15\% moisture & 28 Feb 2027\\
5 & Confirm CH₄ treatment, persistence, buffer and H/Corg threshold with the standard and verifier; re-baseline the model & Carbon / MRV lead & Written confirmation; plan re-issued & 31 Jan 2027\\\bottomrule
\end{tabularx}}

\hd{What I would validate with real data first}
\begin{itemize}
\item Kiln spec and the first month of batch logs: t dry per batch, batches per running kiln-day, utilisation. These dominate the tornado.
\item Rejection reasons and remediation rate from the January 2027 first-issuance package.
\item Moisture at the kiln against CH₄, and lab C\textsubscript{org} and H/C\textsubscript{org} by residue type (cotton, pigeon pea, castor).
\item What ``>100,000 t planned for 2027'' counts: 2027 issuances, which may include late-2026 production, or 2027 vintage.
\item The village payment ledger by site, to test whether supply really lags payment experience by about a month.
\end{itemize}
\vfill
{\scriptsize\color{muted}Independent analysis using public information and synthetic data. Not affiliated with or endorsed by Varaha. Model: 218 parameters in assumptions.yaml (5 public, 15 literature, 198 assumptions); 5,000 Monte Carlo runs; Kon-Tiki values from Schmidt et al.\ 2016 (PLOS ONE). Every figure is regenerated by tools/build\_docs.py.}
\end{document}
